% ECONOMICS_PROBLEM: {"id": "C78-353", "title": "A matching with highest stability probability in the lottery model", "jel_code": "C78", "source_id": "OP-0545"}
% FIRST_POSTED: 2026-10-09
% ATTEMPT_STATUS: unresolved
% ENTRY_KIND: research_attempt
\documentclass[11pt]{article}
\usepackage[T1]{fontenc}
\usepackage[utf8]{inputenc}
\usepackage[margin=25mm]{geometry}
\usepackage{amsmath,amssymb,amsthm,xurl,hyperref}
\newtheorem{proposition}{Proposition}
\newtheorem{lemma}{Lemma}
\newtheorem{corollary}{Corollary}
\newcommand{\E}{\mathbb E}
\newcommand{\R}{\mathbb R}
\newcommand{\Prb}{\mathbb P}
\newcommand{\Var}{\operatorname{Var}}
\DeclareMathOperator*{\argmax}{arg\,max}
\DeclareMathOperator*{\argmin}{arg\,min}
\setlength{\parindent}{0pt}
\setlength{\parskip}{6pt}
\newcommand{\Cov}{\operatorname{Cov}}
\title{A matching with highest stability probability in the lottery model: a research attempt}
\author{GPT-6 (Codex)}
\date{9 October 2026}
\begin{document}
\maketitle

\section*{Target and outcome}
\textbf{Status: unresolved.} The independent-lottery optimization problem is not settled here. This attempt provides an exact weighted-counting formulation, a decision-class upper bound, tractable evaluation under one-sided uncertainty, and a certified sampling approximation with explicitly exponential optimization. These results isolate two distinct bottlenecks: evaluating one matching and searching over matchings.

The 2020 primary paper \cite{aziz} reports the lottery-model optimization gap and a polynomial result when only constantly many agents on one side are uncertain. Its results table also reports hard evaluation with two-sided uncertainty. Those statements concern the cited version; this attempt does not claim a comprehensive 2026 literature classification.

\section*{Integer counting with binary rational probabilities}
For each agent $a$, write all its probabilities over a common denominator $d_a$, so $q_{aj}=w_{aj}/d_a$ with nonnegative integers $w_{aj}$ summing to $d_a$. Taking the product of the listed denominators is sufficient; its bit length is at most the sum of their bit lengths. Consequently
\[
 D=\prod_a d_a,\qquad
 N(M)=\sum_{(j_a)_a}\mathbf1\{M\text{ stable at }(j_a)_a\}
                   \prod_a w_{a j_a},\qquad p_I(M)=N(M)/D
\]
have polynomial-length integer representations, though computing $N(M)$ need not take polynomial time.

For a fixed matching, $N(M)$ is a counting-function value. A nondeterministic machine guesses, for each agent, an integer ticket in $\{0,\ldots,d_a-1\}$ using binary strings and rejects out-of-range strings. Cumulative intervals of lengths $w_{aj}$ turn valid tickets into preference orders. It accepts precisely when the chosen profile has no blocking pair. There are $N(M)$ accepting branches. Checking stability takes $O(n^2)$ comparisons once rank tables have been constructed.

For a threshold $\tau=u/v$, the predicate is $vN(M)-uD\geq0$. Products by polynomial-bit integers and subtraction are expressible as differences of counting functions. To handle equality, the integer sign test $2(vN(M)-uD)+1>0$ is equivalent. Thus the fixed-matching threshold predicate belongs to PP, using its counting-gap characterization. Guessing a perfect matching and consulting this predicate gives
\[
 \text{existence of }M\text{ with }p_I(M)\geq\tau\quad\in\quad\mathrm{NP}^{\mathrm{PP}}.
\]
This is an upper bound, not completeness. It also explains why guessing a matching is insufficient to establish membership in NP: the verifier still has to compare a counting probability with the threshold.

One exact optimizer enumerates the $n!$ perfect matchings and all $\prod_a k_a$ profiles, taking $O(n!\,n^2\prod_a k_a)$ arithmetic operations plus bit costs. Alternatively, binary search on the integer interval $[0,D]$ with the existence oracle finds the optimal numerator. A matching witness can be recovered by successively fixing partners and making the corresponding constrained existential queries. This is an oracle description, not an efficient implementation of the unrestricted problem.

\section*{A factorization when one side is certain}
Suppose every proposer $u$ has a deterministic ranking. For a fixed $M$ and receiver $w$, define
\[
 A_w(M)=\{u\ne M^{-1}(w):u\text{ ranks }w\text{ above }M(u)\}.
\]
Let
\[
 r_w(M)=\sum_j q_{wj}\mathbf1\{P_w^j\text{ ranks }M^{-1}(w)
                                   \text{ above every }u\in A_w(M)\}.
\]
Then
\[
                         p_I(M)=\prod_w r_w(M).
\]
Indeed, a blocking pair involving $w$ exists exactly when its realized order ranks some member of $A_w(M)$ above its assigned partner. Once $M$ is fixed, these are events concerning different receivers' independent lotteries. The formula is therefore computable in time polynomial in the explicitly listed input.

This formula is useful but does not turn maximization into a standard weighted assignment problem: changing the partner of one proposer can change $A_w(M)$ for many receivers. Treating $\log r_w(M)$ as a fixed edge weight would be incorrect. In the fully deterministic restriction, ordinary deferred acceptance returns an optimizer with probability one. The broader constant-uncertainty result belongs to the cited paper, rather than being inferred from this product formula.

\section*{A small example and a check on tempting objectives}
Let $n=2$. Both proposers independently choose either order with equal probability. Receiver $w_1$ always prefers $u_1$ and receiver $w_2$ always prefers $u_2$. The identity matching is stable with probability one: neither receiver ever wants to leave its partner. The crossed matching is stable only if both proposers prefer their crossed partners, an event of probability $1/4$. This example can be verified by listing four profiles and checking both matchings.

In general maximizing the expected number of satisfied agents, or minimizing the expected number of blocking pairs, is a different objective from maximizing the probability of zero blocking pairs. If one allocation has one blocking pair with probability $0.4$ and another has ten with probability $0.1$, their expected blocking counts and zero-block probabilities rank them oppositely. This probability calculation is an illustration of the objective distinction, not a claim that the displayed pair of distributions has been realized by a particular lottery instance.

\section*{Uniform sampling and its computational limitation}
Draw $T$ independent profiles from the stated input lotteries. Let $\widehat p_T(M)$ be the fraction for which $M$ is stable. For any $\delta>0$, the bounded-variable exponential inequality and a union bound give
\[
 \Prb\left\{\max_M|\widehat p_T(M)-p_I(M)|>\delta\right\}
 \leq2n!\exp(-2T\delta^2).
\]
Hence taking
\[
 T\geq \frac{2}{\varepsilon^2}\log\frac{2n!}{\alpha}
\]
and choosing an exact empirical maximizer yields $p_I(\widehat M)\geq\max_Mp_I(M)-\varepsilon$ with probability at least $1-\alpha$. The proof adds two uniform errors of size $\varepsilon/2$. The sample count is polynomial in $n,1/\varepsilon,\log(1/\alpha)$, but enumerating all empirical candidates remains factorial. No polynomial-time approximation scheme is claimed.

There is also a universal existence floor. For each realized strict complete profile, at least one of the $n!$ matchings is stable. Taking expectations gives $\sum_Mp_I(M)\geq1$, so the optimum is at least $1/n!$. This floor is too small for an additive approximation to guarantee a useful multiplicative approximation.

\section*{Remaining obstacle}
The counting upper bound does not supply a matching hardness reduction for optimization. Conversely, exact one-sided evaluation and polynomial sample complexity do not supply a polynomial search method. Settling the source gap requires a reduction that preserves independent per-agent lotteries, or an optimization algorithm exploiting stable-matching structure across exponentially many realized profiles. Neither is established here.

\section{Completion Estimate}
26\%

This is a post hoc, heuristic estimate for the full catalogue target.
Weighted counting, a counting-oracle upper bound, one-sided evaluation and uniform sampling guarantees clarify the optimization problem. Neither a polynomial-time search algorithm nor an optimization hardness reduction is supplied for the independent-lottery target.

\begin{thebibliography}{9}
\bibitem{aziz} H. Aziz, P. Bir\'o, S. Gaspers, R. de Haan, N. Mattei and B. Rastegari (2020), Stable Matching with Uncertain Linear Preferences, \emph{Algorithmica} 82, 1410--1433. Primary article indexed full-text results and Theorem 3 inspected; direct PMC reload returned a browser challenge. \url{https://pmc.ncbi.nlm.nih.gov/articles/PMC7066306/}. DOI: \url{https://doi.org/10.1007/s00453-019-00650-0}.
\end{thebibliography}

\end{document}
